% Part:sets-functions-relations
% Chapter: size-of-sets
% Section: enumerations-alt

\documentclass[../../../include/open-logic-section]{subfiles}

\begin{document}

\olfileid[ar]{sfr}{siz}{enm-alt}

\olsection{التعدادات ومجموعات \usetoken{S}{enumerable}}

\begin{editorial}
التعداد هنا دالة تقابلية مع $\Nat$ أو مع مقطع ابتدائي منها، على اصطلاح نظرية المجموعات. فبين هذا التعريف وتعريفات \olref[enm]{sec} فرق يسير، وترد هنا أمثلة ذلك القسم كلها مرة أخرى. وهذا العرض أوجز منه بعض الإيجاز.
\end{editorial}

من طرق تعيين المجموعة المنتهية أن نسرد !!{element}sها؛ ومن ذلك أن نكتب
\[
  A = \{a_1, a_2, \ldots, a_n\}.
\]
فإذا كانت !!{element}s~$a_1$، \dots،~$a_n$ متمايزة كلها، نشأ من هذا السرد !!a{bijection} بين $A$ والأعداد الطبيعية الأولى التي عددها $n$، وهي $0$، \dots،~$n-1$. وبالعكس، كل مجموعة منتهية يمكن أن نجعلها في مثل هذه المراسلة، لأن عدد !!{element}sها منته. فمتى كانت $A$ منتهية وجدت !!a{bijection} بين $A$ و$\{0, \dots, n-1\}$، حيث $n$ عدد !!{element}s~$A$.

وإذا أجزنا ضروبًا من المجموعات اللانهائية أمكن أن نجيز لبعضها تعدادًا أيضًا. وضبط هذا أن تعداد المجموعة اللانهائية يكون بـ!!a{bijection} بينها وبين مجموعة الأعداد الطبيعية بأسرها~$\Nat$.

\begin{defn}[التعداد في اصطلاح نظرية المجموعات]
\emph{تعداد} المجموعة $A$ هو !!a{bijection} مداها $A$، ومجالها مقطع ابتدائي من الأعداد الطبيعية $\{0, 1, \ldots, n\}$ {أو} مجموعة الأعداد الطبيعية كلها~$\Nat$.
\end{defn}

\begin{explain}
واختيار اسم \emph{التعداد} موافق للتصور؛ فقولنا إننا عددنا المجموعة $A$ معناه وجود !!a{bijection}~$f$ نحصي بها عناصر المجموعة~$A$ واحدًا بعد واحد: ذو الرتبة~$0$ هو $f(0)$، والذي يليه $f(1)$، \ldots وذو الرتبة~$n$ هو $f(n)$\ldots.\footnote{فالعد هنا مبتدئ من $0$؛ ولو ابتدأناه من~$1$ لم يتغير المقصود تغيرًا ذا بال، وإنما نستبدل~$\Nat$ بـ~$\PosInt$.} ويزداد هذا المعنى بيانًا بالتعريف التالي.
\end{explain}

\begin{defn}
\ollabel{defn:enumerable}
المجموعة~$A$ !!{enumerable} إذا وفقط إذا كان $A = \emptyset$ أو وجد لـ~$A$ تعداد. وأما $A$ فتسمى !!{nonenumerable} إذا وفقط إذا لم تكن $A$ !!{enumerable}.
\end{defn}

\begin{explain}
فحاصل الأمر أن المجموعة !!{enumerable} إما خالية، وإما ذات تعداد يتيح عد !!{element}sها واحدًا واحدًا.
\end{explain}

\begin{ex}
يكفينا لتعداد الأعداد الطبيعية دالة الهوية $\Id{\Nat} \colon \Nat \to \Nat$ التي يحدها $\Id{\Nat}(n) = n$. وأما الأعداد الطبيعية \emph{الموجبة}، أي $\Nat^+ =
\Nat \setminus \{0\}$، فتعددها دالة اللاحق $g(n) = n + 1$.
\end{ex}

\begin{prob}
برهن أن $A$ !!{enumerable} إذا وفقط إذا كان $A = \emptyset$ أو وجدت !!a{surjection} $f\colon \Nat \to A$. ثم برهن أن $A$ !!{enumerable} إذا وفقط إذا وجدت !!a{injection} $g\colon A \to \Nat$.
\end{prob}

\begin{ex}
لتكن الدالتان $f\colon \Nat \to \Nat$ و$g \colon \Nat \to \Nat$ معينتين بالحدين
\begin{align*}
  f(n) & = 2n \text{ و}\\
  g(n) & = 2n+1
\end{align*}
فالأولى تعدد الأعداد الطبيعية الزوجية، والثانية تعدد الفردية. لكن لا واحدة منهما !!{surjective} إلى المجال المقابل المعين؛ ولذلك ليست أي منهما تعدادًا لـ$\Nat$.
\end{ex}

\begin{prob}
عين دالة تعدد الأعداد المربعة $1$، $4$، $9$، $16$، \dots
\end{prob}

\begin{ex}
نكتب $\lceil x \rceil$ لـ\emph{دالة السقف}، وقيمتها أصغر عدد صحيح لا يقل عن $x$. فالدالة $f \colon \Nat \to \Int$ المعرفة بالحد
\[
  f(n) = (-1)^{n} \left\lceil\tfrac{n}{2}\right\rceil
\]
تعدد مجموعة الأعداد الصحيحة~$\Int$، كما يظهر من القيم الآتية:
\[
\begin{array}{c c c c c c c c}
f(0) & f(1) & f(2) & f(3) & f(4) & f(5) & f(6) & \dots \\ \\
\big\lceil \tfrac{0}{2} \big\rceil & -\big\lceil \tfrac{1}{2}\big\rceil &  \big\lceil \tfrac{2}{2} \big\rceil & -\big\lceil \tfrac{3}{2} \big\rceil & \big\lceil \tfrac{4}{2} \big\rceil  & -\big\lceil \tfrac{5}{2}\big\rceil & \big\lceil \tfrac{6}{2} \big\rceil & \dots \\ \\
0 & -1 & 1 & -2 & 2 & -3 & 3& \dots
\end{array}
\]
وترى أن $f$ تستوفي قيم $\Int$ بالتنقل بين الأعداد الصحيحة الموجبة والسالبة. ويجوز أن نعين $f$ بحالتين على الوجه الآتي:
\[
f(n) = \begin{cases}
  \frac{n}{2} & \text{إذا كان $n$ زوجيًا}\\
  -\frac{n+1}{2} & \text{إذا كان $n$ فرديًا}
  \end{cases}
\]
\end{ex}

\begin{prob}
إذا كانت $A$ و$B$ !!{enumerable}، فبين أن $A \cup B$ كذلك.
\end{prob}

\begin{prob}
استعمل الاستقراء على $n$ لإثبات أن $A_1$، $A_2$، \dots،~$A_n$ إذا كانت كلها !!{enumerable}، كان $A_1 \cup \dots \cup A_n$ كذلك.
\end{prob}

\end{document}
